% =====================================================================
%  resumen.tex: Resumen y palabras clave. Se imprime en la hoja
%  «Documento Técnico» (frontmatter/documento-tecnico.tex).
% =====================================================================
El presente Trabajo Terminal desarrolla Qadra, un prototipo de sistema web
para la gestión de expedientes digitales y la actividad procesal de un
despacho legal en materia penal, con servicios criptográficos de
\emph{integridad} y \emph{no repudio}. El servidor, escrito en Rust con
arquitectura hexagonal, administra expedientes conforme al Código Nacional
de Procedimientos Penales: etapas, audiencias, medidas cautelares,
resoluciones, recursos y plazos calculados con calendarios jurisdiccionales.
Cada versión documental se cifra con AES-256-GCM bajo una llave propia y
cada operación queda en una bitácora encadenada que detecta alteraciones.
La firma digital RSA-3072 se produce en el equipo del titular, sin que la
clave privada salga de él, y una autoridad de sellado de tiempo local
emite sellos RFC~3161; con ambos se firman documentos individuales y
contratos entre el despacho y su cliente. El acceso combina contraseñas
protegidas con Argon2id, un segundo factor TOTP y control por roles. Cada
firma se exporta como un paquete de evidencia que puede verificarse con
OpenSSL o en el navegador, sin depender del sistema. El prototipo está
desplegado; su ejecución de cierre aprobó 6\,397 pruebas automatizadas del
servidor, además de las pruebas de la interfaz y de navegador, y la
usabilidad se evaluó con dos personas ajenas al despacho.

\vspace{1em}
\noindent\textbf{Palabras clave:} expedientes digitales, integridad,
no repudio, firma digital, criptografía aplicada, gestión legal, CNPP.
